% =====================================================================
% Приложение В. Состав программного репозитория.
% Параграфный формат для компактности.
% =====================================================================

\section{Состав программного репозитория}
\label{app:repo}

В приложении приведено краткое описание состава программного
репозитория проекта. Состав фиксирован по состоянию идентификатора
фиксации, указанного в приложении~\ref{app:github}.

\begingroup
\small

\subsection*{Корневая структура}

\noindent\textbf{\code{src/atm/}} --- основной пакет фреймворка,
тринадцать функциональных слоев.\quad
\textbf{\code{conf/}} --- декларативная конфигурация экспериментов
и компонентов в формате YAML.\quad
\textbf{\code{scripts/}} --- самостоятельные скрипты обслуживания
экспериментов и сборки производных артефактов.\quad
\textbf{\code{notebooks/}} --- шаблон Jupyter-ноутбука для
аналитической постобработки результатов прогонов.\quad
\textbf{\code{pyproject.toml}} --- декларация пакета, зависимостей
и точки входа \code{atm}.\quad
\textbf{\code{README.md}} --- точка входа в документацию проекта
и инструкция по сборке среды.

\subsection*{Слои пакета \code{src/atm/}}

\noindent\textbf{\code{atm.core}} --- контейнеры состояния и
базовые типы (\code{AgentState}, \code{SharedState},
\code{GraphState}, \code{Phase}, \code{HumanRole}).\quad
\textbf{\code{atm.llm}} --- оболочка над провайдерами языковых
моделей, ценообразование, бюджетные ограничители,
детерминированный \code{FakeLLM}, фабрика \code{build\_llm}.\quad
\textbf{\code{atm.tools}} --- реестр инструментов и обертка
над Docker-песочницей.\quad
\textbf{\code{atm.agents}} --- базовый класс \code{Agent} и
пять специализированных ролей плюс координатор.\quad
\textbf{\code{atm.topology}} --- шесть реализаций топологий
(\code{star.py}, \code{chain.py}, \code{mesh.py}, \code{debate.py},
\code{hierarchical.py}, \code{adaptive.py}) поверх общего протокола
\code{Topology}.\quad
\textbf{\code{atm.phases}} --- конечный автомат фаз, маршрутизатор
топологий, защитные правила переключения, ворота передачи
состояния.\quad
\textbf{\code{atm.human}} --- протокол \code{HumanGateway},
реализации шлюзов, маршрутизатор ролей, политика тайм-аутов.\quad
\textbf{\code{atm.storage}} --- объектно-реляционное отображение,
модели таблиц PostgreSQL, асинхронные сессии, писатель Parquet,
обертка над чекпойнтером LangGraph, миграции \code{alembic/}.\quad
\textbf{\code{atm.observability}} --- перехватчик событий
LangGraph, сериализация, обертки над логгером.\quad
\textbf{\code{atm.tasks}} --- загрузчики и оценщики четырех
корпусов (HumanEval, GSM8K, CommonGen, InfiAgent-DABench).\quad
\textbf{\code{atm.evaluation}} --- судья оценки качества, оракулы
достоверности, агрегатор шкалы NASA-TLX.\quad
\textbf{\code{atm.experiment}} --- схемы конфигов на Pydantic,
загрузчик OmegaConf с расширением по сетке, командная строка
\code{atm} на Typer.\quad
\textbf{\code{atm.analysis}} --- загрузчики результатов,
агрегаторы метрик, генератор графиков, построитель оракула.

\subsection*{Конфигурационные файлы \code{conf/}}

\noindent\textbf{\code{conf/experiments/}} --- полные конфигурации
экспериментов E1--E4, подтверждающих прогонов и проверок
работоспособности.\quad
\textbf{\code{conf/agents/}} --- системные подсказки шести ролей
агентов.\quad
\textbf{\code{conf/topology/}} --- параметры по умолчанию для
шести топологий.\quad
\textbf{\code{conf/human/}} --- параметры шлюза HITL и системные
подсказки симулятора.\quad
\textbf{\code{conf/model/}} --- параметры языковых моделей.\quad
\textbf{\code{conf/task/}, \code{conf/tasks/}} --- загрузчики
корпусов и стратифицированные выборки.\quad
\textbf{\code{conf/oracle/}} --- конфигурации построения оракула.\quad
\textbf{\code{conf/evaluation/}} --- параметры судьи и агрегаторов.\quad
\textbf{\code{conf/sandbox/}} --- параметры Docker-песочницы и
профиль системных вызовов.

\subsection*{Скрипты и ноутбуки}

\noindent\textbf{\code{scripts/gen\_analysis\_notebook.py}} ---
генератор аналитического ноутбука из шаблона.\quad
\textbf{\code{scripts/derive\_seccomp.py}} --- сборка профиля
системных вызовов для песочницы.\quad
\textbf{\code{notebooks/analysis\_template.ipynb}} --- шаблонный
ноутбук из шести секций для постобработки и визуализации
результатов прогона.

\endgroup
